\documentclass[11pt,a4paper]{article}

\usepackage{amsmath, amssymb, amsthm}
\usepackage{physics}
\usepackage{bm}
\usepackage{siunitx}
\usepackage{booktabs}
\usepackage{array}
\usepackage[margin=2cm]{geometry}
\usepackage{enumitem}
\usepackage{hyperref}

\renewcommand{\arraystretch}{1.6}

\title{Technical Note: Multi-Sphere T-Matrix Theory\\[4pt]
\large Mathematical Formulation and Computational Implementation}
\author{N.\ Moteki}
\date{Last updated: 2026-11-20 (MSTMforCAS.jl v1.1.1)}

\begin{document}
\maketitle

\begin{abstract}
This document provides a self-contained description of the mathematical theory behind the Multi-Sphere T-Matrix (MSTM) method as implemented in MSTMforCAS.jl.
It covers the VSWF expansion formalism, Lorenz--Mie theory for single spheres, the multi-sphere interaction equation, VSWF translation addition theorem, the Complex BiConjugate Gradient (CBICG) iterative solver, the FFT-accelerated translation algorithm, and the output quantities (scattering amplitudes and cross-section efficiencies).
The implementation follows the formulation of Mackowski (1991, 1994, 1996), Xu (1996), Mackowski \& Mishchenko (2011), and Mackowski \& Kolokolova (2022).
Scattering amplitude conventions follow Bohren \& Huffman (1983, hereafter BH83).
\end{abstract}

\tableofcontents

%% ============================================================
\section{Problem Setup}
\label{sec:problem-setup}

\subsection{Geometry and incident field}

Consider $N_S$ homogeneous, non-overlapping spheres embedded in an infinite, lossless, homogeneous medium of real refractive index $n_\mathrm{med}$.
Sphere $i$ ($i = 1, \dots, N_S$) has centre position $\bm{r}_i$, radius $a_i$, and complex refractive index $m_\mathrm{sphere}$ (common to all spheres in the parameter-sweep mode of this code).
A time-harmonic plane wave ($e^{-i\omega t}$ convention, following BH83) propagates along the $+z$ axis in the medium with vacuum wavelength $\lambda_0$ and wavenumber in the medium
\begin{equation}
  k = \frac{2\pi \, n_\mathrm{med}}{\lambda_0}.
  \label{eq:wavenumber}
\end{equation}

A positive imaginary part of $m_\mathrm{sphere}$ corresponds to absorption.
All lengths are non-dimensionalised by $k$: positions become $k\bm{r}_i$, radii become size parameters $x_i = k a_i$, and the relative refractive index entering all scattering formulas is $m_\mathrm{rel} = m_\mathrm{sphere} / n_\mathrm{med}$.
Tildes are suppressed hereafter; all equations are in dimensionless units unless stated otherwise.

The volume-equivalent sphere radius $a_\mathrm{eff}$ and its size parameter $x_\mathrm{eff}$ are defined as
\begin{equation}
  a_\mathrm{eff} = \left(\sum_{i=1}^{N_S} a_i^3\right)^{\!1/3}, \qquad
  x_\mathrm{eff} = k\,a_\mathrm{eff}.
  \label{eq:a-eff}
\end{equation}
These appear in the normalisation of cross-section efficiencies (Sec.~\ref{sec:output}).

\subsection{Vector spherical wave function expansion}

The electromagnetic field in each region is expanded in vector spherical wave functions (VSWFs) $\bm{M}_{mn}^{(\nu)}$, $\bm{N}_{mn}^{(\nu)}$ of degree $m$ and order $n$ ($n = 1, 2, \dots$; $m = -n, \dots, n$), where $\nu = 1$ denotes regular (Bessel $j_n$) waves and $\nu = 3$ denotes outgoing (Hankel $h_n^{(1)}$) waves.
The normalisation convention for these VSWFs follows that of Mishchenko et al.\ (2002) and Mackowski (1994).
The scattered field from sphere $i$ is expanded in outgoing VSWFs centred at $\bm{r}_i$:
\begin{equation}
  \bm{E}_\mathrm{sca}^{(i)}(\bm{r}) = \sum_{n=1}^{N_i} \sum_{m=-n}^{n}
  \left[ a_{mn}^{(i)}\,\bm{M}_{mn}^{(3)}(\bm{r} - \bm{r}_i)
       + b_{mn}^{(i)}\,\bm{N}_{mn}^{(3)}(\bm{r} - \bm{r}_i) \right]
  \label{eq:sca-expansion}
\end{equation}
where $N_i$ is the truncation order for sphere $i$.
The field exciting sphere $i$ --- the incident plane wave plus the scattered fields from all other spheres --- is likewise expanded in regular VSWFs $\bm{M}_{mn}^{(1)}, \bm{N}_{mn}^{(1)}$ centred at $\bm{r}_i$.

\subsection{Left--right circular (lr\_tran) basis}

Following Mackowski (1996, \S 2), the code works in the left--right circular VSWF basis obtained by the linear combinations
\begin{equation}
  \tilde{\bm{N}}_{mn,L}^{(\nu)} = \bm{N}_{mn,1}^{(\nu)} + \bm{N}_{mn,2}^{(\nu)}, \qquad
  \tilde{\bm{N}}_{mn,R}^{(\nu)} = \bm{N}_{mn,1}^{(\nu)} - \bm{N}_{mn,2}^{(\nu)}
  \label{eq:lr-basis}
\end{equation}
where the subscripts 1, 2 refer to the TE and TM modes.
The fundamental mathematical advantage of this basis is that it completely decouples the translation matrix $\mathbf{H}_{ij}$ between the two polarisation blocks.
While the single-sphere T-matrix is diagonal in the combined multipole index $mn = n(n+1) + m$ in both the standard TE/TM and lr\_tran bases, in the lr\_tran basis it couples the two polarisation blocks (labeled $p = 1$ for left and $p = 2$ for right).
This structure greatly simplifies the multi-sphere interaction equations, as the computationally expensive translation operations can be performed independently for each polarisation block.

%% ============================================================
\section{Single-Sphere Scattering: Lorenz--Mie Theory}
\label{sec:mie}

For a homogeneous sphere with size parameter $x$ and relative refractive index $m_\mathrm{rel}$, the Lorenz--Mie scattering coefficients $a_n$ (TM/electric) and $b_n$ (TE/magnetic) are given by (BH83, Eq.~4.88):
\begin{equation}
  a_n = \frac{\left[\dfrac{D_n(m_\mathrm{rel}\,x)}{m_\mathrm{rel}} + \dfrac{n}{x}\right] \psi_n(x) - \psi_{n-1}(x)}
             {\left[\dfrac{D_n(m_\mathrm{rel}\,x)}{m_\mathrm{rel}} + \dfrac{n}{x}\right] \xi_n(x) - \xi_{n-1}(x)}
  \label{eq:mie-an}
\end{equation}
\begin{equation}
  b_n = \frac{\left[m_\mathrm{rel}\,D_n(m_\mathrm{rel}\,x) + \dfrac{n}{x}\right] \psi_n(x) - \psi_{n-1}(x)}
             {\left[m_\mathrm{rel}\,D_n(m_\mathrm{rel}\,x) + \dfrac{n}{x}\right] \xi_n(x) - \xi_{n-1}(x)}
  \label{eq:mie-bn}
\end{equation}
where:
\begin{itemize}[nosep]
  \item $\psi_n(z) = z\,j_n(z)$ is the Riccati--Bessel function ($j_n$: spherical Bessel function of the first kind)
  \item $\xi_n(z) = z\,h_n^{(1)}(z) = \psi_n(z) + i\,\chi_n(z)$ is the Riccati--Hankel function
  \item $\chi_n(z) = -z\,y_n(z)$ ($y_n$: spherical Bessel function of the second kind)
  \item $D_n(z) = \frac{d}{dz}\ln[z\,j_n(z)] = \psi_n'(z)/\psi_n(z)$ is the logarithmic derivative
\end{itemize}

The logarithmic derivative $D_n(y)$ with $y = m_\mathrm{rel}\,x$ is computed by downward recurrence (Wiscombe 1980):
\begin{equation}
  D_{n-1}(y) = \frac{n}{y} - \frac{1}{D_n(y) + n/y},
  \label{eq:dn-recurrence}
\end{equation}
starting from a sufficiently large $n_\mathrm{start} \geq \max(N_\mathrm{max}, |y|) + 16$ where $D_{n_\mathrm{start}} \approx 0$.
This recurrence is unconditionally stable for any complex $y$.
The Riccati--Bessel functions $\psi_n(x)$ and $\chi_n(x)$ (for real argument $x$) are computed by upward recurrence:
\begin{equation}
  f_n(x) = \frac{2n - 1}{x}\,f_{n-1}(x) - f_{n-2}(x)
  \label{eq:bessel-recurrence}
\end{equation}
with starting values $\psi_0 = \sin x$, $\psi_1 = \sin x / x - \cos x$, $\chi_0 = -\cos x$, $\chi_1 = -\cos x / x - \sin x$.

The multipole expansion is truncated at order $N_i$ determined by the criterion of Wiscombe (1980) as implemented in the MSTM code:
\begin{equation}
  N_i = \mathrm{round}\!\left(x_i + 4\,x_i^{1/3}\right) + 5,
  \label{eq:nstop}
\end{equation}
refined by a convergence check: starting from this upper bound, $N_i$ is reduced until the cumulative partial extinction efficiency $Q_\mathrm{ext}^{(N_i)}$ deviates from the total by more than $10^{-6}$.
This criterion describes an isolated sphere.
For spheres that touch, the order can be raised by the two-sphere criterion of Sec.~\ref{sec:contact-order}.

In the lr\_tran basis, the single-sphere T-matrix for multipole order $n$ takes the $2 \times 2$ form coupling the $p = 1$ and $p = 2$ blocks (Mackowski 1996, derived from the \texttt{mieoa} routine in MSTM v4.0):
\begin{equation}
  \mathbf{T}_n = -\frac{1}{2}
  \begin{pmatrix} a_n + b_n & a_n - b_n \\ a_n - b_n & a_n + b_n \end{pmatrix}
  \label{eq:tmatrix-lr}
\end{equation}
This matrix is diagonal in the combined index $mn$ (it does not mix different $(m,n)$ pairs) but couples the two circular polarisation blocks.
For an isotropic sphere, $a_n \neq b_n$ in general, so the off-diagonal coupling is physically meaningful and arises from the difference in electric and magnetic multipole scattering.

%% ============================================================
\section{Multi-Sphere Interaction Equation}
\label{sec:multi-sphere}

\subsection{Coupled field equations}

The field that excites sphere $i$ consists of the external incident wave plus the scattered fields from all other spheres, translated to the coordinate origin of sphere $i$ using the VSWF translation addition theorem (Mackowski 1991, 1996; Xu 1996):
\begin{equation}
  \bm{p}_\mathrm{exc}^{(i)} = \bm{p}_\mathrm{inc}^{(i)} + \sum_{j \neq i} \mathbf{H}_{ij}\,\bm{a}^{(j)}
  \label{eq:exciting-field}
\end{equation}
where $\mathbf{H}_{ij}$ is the translation matrix that re-expands outgoing VSWFs centred at $\bm{r}_j$ into regular VSWFs centred at $\bm{r}_i$ (Hankel-to-Bessel type), and $\bm{a}^{(j)}$ is the vector of scattered-field expansion coefficients of sphere $j$.

The T-matrix relates the exciting field to the scattered field:
\begin{equation}
  \bm{a}^{(i)} = \mathbf{T}_i\,\bm{p}_\mathrm{exc}^{(i)}.
  \label{eq:tmatrix-relation}
\end{equation}

Substituting and rearranging into a single global linear system for $\bm{x} = (\bm{a}^{(1)}, \dots, \bm{a}^{(N_S)})^\top$ gives (Mackowski 1994, Eq.~6):
\begin{equation}
  (\mathbf{I} - \mathbf{T}\,\mathbf{A})\,\bm{x} = \mathbf{T}\,\bm{p}_\mathrm{inc}
  \label{eq:global-system}
\end{equation}
where $\mathbf{T} = \mathrm{diag}(\mathbf{T}_1, \dots, \mathbf{T}_{N_S})$ is the block-diagonal T-matrix and $\mathbf{A}$ is the off-diagonal translation operator with blocks $\mathbf{H}_{ij}$ ($i \neq j$) and zeros on the diagonal.
The dimension of this system is $\sum_i 2\,N_i(N_i+2)$ --- twice the multipole block size because each sphere carries both $p = 1$ and $p = 2$ polarisation blocks.

\subsection{Truncation order for touching spheres}
\label{sec:contact-order}

The single-sphere criterion of Eq.~\eqref{eq:nstop} resolves the field scattered by an isolated sphere.
It does not resolve the near field between two spheres that almost touch.
For small spheres the order needed for this coupling is set by $m_\mathrm{rel}$ and by the contact geometry.
It is nearly independent of the size parameter (Mackowski 1994).
For small, strongly absorbing monomers the forward amplitudes of an aggregate can be off by several percent at the single-sphere order.

The code can choose the order from a preliminary two-sphere calculation (option \texttt{truncation\_eps}, function \texttt{contact\_truncation\_order}).
Two equal spheres of size parameter $x$ are placed on the $x$ axis, perpendicular to the incident direction.
One incident polarisation then has its electric field along the axis of the pair, which is the slowly converging configuration.
The centre distance is
\begin{equation}
  d = 2x\,(1 + g),
  \label{eq:contact-distance}
\end{equation}
where $g$ is the surface separation divided by the sum of the radii.
The default $g = 0$ (point contact) is the slowest case.
The forward amplitudes $S_1(0^\circ)$ and $S_2(0^\circ)$ of the pair are computed with a common order $N$ for $N = N_\mathrm{Mie}, N_\mathrm{Mie}+1, \dots$, where $N_\mathrm{Mie}$ is the order of Eq.~\eqref{eq:nstop}.
The relative increment at order $N$ is
\begin{equation}
  \delta_N = \frac{\left[\,\sum_{k=1}^{2} \left|S_k^{(N)} - S_k^{(N-1)}\right|^2\right]^{1/2}}
                  {\left[\,\sum_{k=1}^{2} \left|S_k^{(N)}\right|^2\right]^{1/2}}.
  \label{eq:contact-increment}
\end{equation}
The contribution of the orders above $N$ is bounded under the assumption that the increments decay at least algebraically, $\delta_N \propto N^{-p}$.
The exponent is taken from the last three increments as the smaller of the two estimates
\begin{equation}
  p = \min\!\left[\frac{\ln(\delta_{N-1}/\delta_N)}{\ln\!\left(N/(N-1)\right)},\;
                 \frac{\ln(\delta_{N-2}/\delta_{N-1})}{\ln\!\left((N-1)/(N-2)\right)}\right],
  \label{eq:contact-exponent}
\end{equation}
and the tail bound is
\begin{equation}
  \tau_N = \frac{\delta_N\,N}{p - 1}\;
           \frac{\left[\,\sum_k |S_k^{(N)}|^2\right]^{1/2}}{\min_k |S_k^{(N)}|}
  \qquad (p > 1,\; \delta_N < \delta_{N-1} < \delta_{N-2}).
  \label{eq:contact-tail}
\end{equation}
The bound is taken as infinite if the increments do not yet decrease monotonically or if $p \leq 1$.
The scan stops at the order $N_\mathrm{top}$ at which $\tau_N \leq f\,\epsilon$, where $\epsilon$ is the requested relative accuracy and $f = 0.25$ by default.
It also stops at a cap of $N = 40$, in which case the result is a lower bound and the code issues a warning.

The estimated truncation error at order $N$ is the difference to the highest computed order plus the tail bound,
\begin{equation}
  e_N = \max_{k=1,2} \frac{\left|S_k^{(N)} - S_k^{(N_\mathrm{top})}\right|}{\left|S_k^{(N_\mathrm{top})}\right|} + \tau_{N_\mathrm{top}}.
  \label{eq:contact-error}
\end{equation}
The difference is taken between complex numbers, so $e_N$ covers both amplitude and phase.
The contact order $N_\mathrm{c}$ is the smallest order from which $e_N$ stays below $\epsilon$.
It is never smaller than $N_\mathrm{Mie}$.

In an aggregate calculation every sphere then uses the largest of $N_\mathrm{c}$, its own order from Eq.~\eqref{eq:nstop} and a user-specified order.
In a parameter sweep $N_\mathrm{c}$ is determined for every refractive index.
On the GPU path all refractive indices of a batch share the largest order of the batch.
Suggested values are $\epsilon = 10^{-2}$ for routine use and $\epsilon = 10^{-3}$ for precise calculations.

\subsection{Translation matrix via rotation--axial translation--rotation}

The translation matrix $\mathbf{H}_{ij}$ is computed via the factorisation into rotation, axial translation, and inverse rotation (Mackowski 1996, \S 3; Xu 1996, \S 4), which avoids the direct evaluation of Gaunt integrals and is more efficient:

\begin{enumerate}[nosep]
  \item \textbf{Rotation.}
  Compute the spherical coordinates $(d_{ij}, \theta_{ij}, \phi_{ij})$ of the translation vector $\bm{r}_i - \bm{r}_j$.
  The azimuthal phase factors $e^{im\phi_{ij}}$ and the normalised associated Legendre functions $P_n^m(\cos\theta_{ij})$ provide the rotation from the laboratory frame to a frame where the translation vector is along $z$.

  \item \textbf{Axial translation coefficients.}
  In the rotated frame where the two origins are separated along $z$ by distance $d_{ij}$, the translation coefficients factorise.
  The key quantity involves a sum over intermediate order $w$:
  \begin{equation}
    C_{mn,kl} = \sum_{w=w_\mathrm{min}}^{n+l} c_w(m,n,k,l) \; h_w^{(1)}(d_{ij}) \; P_w^{m-k}(\cos\theta_{ij})
    \label{eq:axial-translation}
  \end{equation}
  where $w_\mathrm{min} = \max(|n-l|, |m-k|)$ and $c_w(m,n,k,l)$ are products of vector coupling (Gaunt) coefficients.
  These coupling coefficients are related to Clebsch--Gordan coefficients and are evaluated using a combined downward--upward recurrence algorithm (Xu 1996, \S 3; Mackowski 1996, Appendix) that avoids explicit factorial computation and maintains numerical stability even at high multipole orders.

  \item \textbf{Polarisation coupling in the lr\_tran basis.}
  In this basis the translation matrix has a particularly simple structure: the even-$w$ and odd-$w$ contributions (relative to $n + l$) combine to give the two polarisation blocks independently.
  Specifically, denoting the even- and odd-$w$ partial sums as $A$ and $B$:
  \begin{equation}
    H_{kl,mn}^{(p=1)} = e^{i(m-k)\phi_{ij}} \left(A + iB\right), \qquad
    H_{kl,mn}^{(p=2)} = e^{i(m-k)\phi_{ij}} \left(A - iB\right).
    \label{eq:lr-translation}
  \end{equation}
  This decoupling is the fundamental advantage of the lr\_tran basis (Mackowski 1996, Eq.~18): the two circular polarisation modes interact with each other only through the T-matrix, not through translation.
\end{enumerate}

The coupling coefficients are precomputed once and cached as a 3D array \texttt{tran\_coef[mn, kl, w]} (Fortran routine \texttt{gentrancoefconstants}), while the distance-dependent Hankel functions and Legendre functions are evaluated per sphere pair.

\subsection{FFT-accelerated translation}

For large numbers of spheres ($N_S \gtrsim 100$), the direct pairwise evaluation of $\mathbf{A}\,\bm{v}$ scales as $O(N_S^2)$ per iteration and dominates the computation time.
Mackowski \& Kolokolova (2022) introduced a discrete Fourier convolution (DFC) scheme into MSTM v4.0 that reduces this to approximately $O(N_S \ln N_S)$ scaling.
The algorithm proceeds as follows:

\begin{enumerate}[nosep]
  \item \textbf{Grid construction.}
  A uniform cubic grid with cell spacing $d$ is mapped onto the target volume.
  The cell spacing is chosen from the mean monomer radius and a target volume fraction $f_v$ as $d = \bar{a} \left(4\pi / 3 f_v\right)^{1/3}$, and each grid dimension is rounded up to the nearest integer of the form $2^a \times 3^b \times 5^c$ for efficient FFT computation (Temperton's generalised prime-factor FFT).

  \item \textbf{Sphere--node association.}
  Each sphere is assigned to the grid cell containing its centre.
  A \emph{neighbour hole} region is defined around each cell: for all cell pairs within this exclusion region (typically $|i - j| \leq \sqrt{2}$, i.e., self + face-adjacent + edge-adjacent cells), sphere--sphere interactions are computed by direct Hankel-type translation, preserving accuracy at short inter-sphere distances.

  \item \textbf{Sphere $\to$ Node.}
  For each sphere, the scattered-field expansion coefficients are translated to its associated grid node using a regular (Bessel $j_n$) translation matrix $\mathbf{J}$.
  This is valid because the sphere centre is within the grid cell (distance less than $d/2$), well within the convergence region of the regular addition theorem.

  \item \textbf{Node-to-node convolution.}
  The grid-based translation operator is a discrete convolution: the translation matrix from node $\bm{i}$ to node $\bm{j}$ depends only on the displacement $\bm{j} - \bm{i}$.
  This convolution is evaluated in the frequency domain:
  \begin{equation}
    \hat{g}_\mu(\bm{k}) = \sum_\nu \hat{H}_{\mu\nu}(\bm{k})\,\hat{a}_\nu(\bm{k})
    \label{eq:fft-convolution}
  \end{equation}
  where $\mu$ and $\nu$ represent the combined multipole indices.
  $\hat{H}_{\mu\nu}(\bm{k})$ is the pre-computed 3D FFT of the cell-to-cell translation matrix (excluding the neighbour-hole cells), $\hat{a}_\nu(\bm{k})$ is the FFT of the zero-padded node coefficient field, and the product is taken element-wise per index pair $(\mu, \nu)$.
  The use of a doubled grid ($2N_x \times 2N_y \times 2N_z$) with zero-padding avoids circular convolution artefacts.

  \item \textbf{Node $\to$ Sphere.}
  The resulting node-level exciting field is translated back to each sphere position via the reverse regular translation.

  \item \textbf{Near-field correction.}
  Direct sphere-to-sphere Hankel-type translation is applied for all sphere pairs within the neighbour-hole region (step 2), adding these contributions to the FFT-based far-field result.
\end{enumerate}

The total cost per iteration becomes $O(N_S + M \log M)$ where $M$ is the number of active grid cells, compared with $O(N_S^2 N_\mathrm{max}^4)$ for the direct algorithm.
The memory overhead consists of the pre-FFT'd translation matrices, which scale as $O(M \times N_\mathrm{node}^2)$ where $N_\mathrm{node}$ is the node multipole order.
Mackowski \& Kolokolova (2022) report speed improvements exceeding two orders of magnitude for systems of $10^3$--$10^4$ spheres.

\paragraph{Far order of the grid.}
\label{sec:fft-far-order}
The memory and the time of the grid part grow steeply with $N_\mathrm{node}$.
A high sphere order, such as the contact order of Sec.~\ref{sec:contact-order}, is needed only for the coupling between spheres that are close to each other.
The code can therefore limit the order that the grid carries (option \texttt{fft\_far\_order}).
Let $F_i \leq N_i$ be the far order of sphere $i$.
The node order is $N_\mathrm{node} = \max\!\left(\lceil d \rceil, \max_i F_i\right)$.
Spheres $i$ and $j$ are called a near pair if
\begin{equation}
  |\bm{r}_i - \bm{r}_j| \leq \kappa\,(a_i + a_j),
  \label{eq:near-pair}
\end{equation}
with $\kappa = 1.5$ by default (option \texttt{fft\_near\_factor}).
The translation operator is split by sphere pair and by multipole order:
\begin{itemize}[nosep]
  \item \emph{Orders up to the far order.}
  Pairs in neighbouring cells are computed directly and all other pairs go through the grid.
  This part equals the operator described above with $N_i$ replaced by $F_i$.
  \item \emph{Orders above the far order.}
  Near pairs are computed directly, whether or not they lie in neighbouring cells.
  All other pairs are neglected.
\end{itemize}
In terms of the blocks $\mathbf{H}_{ij}$, a near pair in neighbouring cells uses the full block.
A pair in neighbouring cells that is not a near pair uses the block truncated to the far orders.
A near pair outside the neighbouring cells uses the full block with the part up to the far orders removed, because the grid already carries that part.
With $F_i = N_i$ the operator is the original one and $\kappa$ has no effect.

In a parameter sweep the far order is $F = \min\!\left[N_\mathrm{max}, \max(F_\mathrm{user}, N_\mathrm{Mie,max})\right]$, where $N_\mathrm{max}$ is the largest sphere order and $N_\mathrm{Mie,max}$ is the largest single-sphere order of Eq.~\eqref{eq:nstop} over the refractive index grid.
The memory and the time of the grid part then do not depend on the orders that the contact criterion adds.
The result has the accuracy of the FFT translation at order $F$.
For an aggregate of 350 monomers at sphere order 8, $F = 4$ reduced the memory of the translation data from \SI{2.7}{\giga\byte} to \SI{0.6}{\giga\byte}.
The forward amplitudes then differed from the direct solution by \SI{0.16}{\percent}, against \SI{0.09}{\percent} with the grid carrying all orders.

%% ============================================================
\section{Iterative Solution: Complex BiConjugate Gradient}
\label{sec:cbicg}

The linear system $\mathbf{L}\,\bm{x} = \bm{b}$ where $\mathbf{L} = \mathbf{I} - \mathbf{T}\,\mathbf{A}$ and $\bm{b} = \mathbf{T}\,\bm{p}_\mathrm{inc}$ is solved by the Complex BiConjugate Gradient (CBICG) method (Mackowski \& Mishchenko 2011, \S 2.3).

\subsection{Operator definitions}

The CBICG iteration requires both the forward operator $\mathbf{L}\,\bm{v} = \bm{v} - \mathbf{T}\,\mathbf{A}\,\bm{v}$ and the adjoint (conjugate transpose) operator $\mathbf{L}^\dagger\,\bm{v} = \bm{v} - \mathbf{A}^\dagger\,\mathbf{T}^\dagger\,\bm{v}$.
The forward operator is computed by applying the translation $\mathbf{A}$ (direct or FFT-accelerated), followed by the block-diagonal T-matrix application $\mathbf{T}$.

The adjoint of the translation operator $\mathbf{A}$ is the computationally expensive part.
Mackowski \& Mishchenko (2011) exploit the S-parity symmetry of the VSWF translation matrix to avoid explicitly forming $\mathbf{A}^\dagger$:
\begin{equation}
  \mathbf{A}^\dagger = \left(\mathbf{S}\,\mathbf{A}\,\mathbf{S}\right)^*
  \label{eq:adjoint-symmetry}
\end{equation}
where $\mathbf{S}$ is the parity operator acting on each sphere's coefficient block as $a_{m,n,p} \mapsto (-1)^m\,a_{-m,n,p}$, and ${}^*$ denotes element-wise complex conjugation.
This identity arises from the spatial inversion symmetry and reciprocity of the electromagnetic translation operator.
The adjoint translation is evaluated by: (i)~applying $\mathbf{S}$, (ii)~performing the same forward translation $\mathbf{A}$, (iii)~applying $\mathbf{S}$ again, and (iv)~complex-conjugating.
This reuses the forward translation code path entirely, including the FFT-accelerated version, at no additional implementation cost.

The system is solved independently for two orthogonal incident linear polarisations ($q = 1$ corresponding to $x$-polarisation, and $q = 2$ corresponding to $y$-polarisation).
The default convergence tolerance is $\varepsilon = 10^{-6}$.

\subsection{CBICG algorithm (default solver)}

The CBICG (Complex BiConjugate Gradient) algorithm is summarised in Table~\ref{tab:cbicg}.  Each iteration requires one forward operator application $\mathbf{L}\,\bm{v}$ and one adjoint application $\mathbf{L}^\dagger\,\bm{v}$, totalling 2 translation operator ($\mathbf{A}$) evaluations per iteration.

\begin{table}[htbp]
\caption{CBICG algorithm (Mackowski \& Mishchenko 2011).}
\label{tab:cbicg}
\centering
\renewcommand{\arraystretch}{1.3}
\begin{tabular}{@{}p{0.95\textwidth}@{}}
\toprule
\textbf{Input:} $\mathbf{L}$, $\bm{b} = \mathbf{T}\,\bm{p}_\mathrm{inc}$, tolerance $\varepsilon$, max iterations $n_\mathrm{max}$. \\
\textbf{Initialise:} $\bm{x}_0 = \bm{b}$, \; $\bm{r}_0 = \bm{b} - \mathbf{L}\,\bm{x}_0$, \; $\bm{q}_0 = \overline{\bm{r}_0}$, \; $\bm{p}_0 = \bm{r}_0$, \; $\bm{w}_0 = \bm{q}_0$, \; $s_0 = \bm{q}_0^\mathsf{H} \bm{r}_0$. \\[4pt]
\textbf{For} $n = 0, 1, \ldots, n_\mathrm{max} - 1$: \\[2pt]
\quad 1. \; $\bm{u} = \mathbf{L}\,\bm{p}_n$ \hfill (1 forward operator application) \\
\quad 2. \; $\bm{v} = \mathbf{L}^\dagger\,\bm{w}_n$ \hfill (1 adjoint operator application via S-parity) \\
\quad 3. \; $\alpha_n = s_n \;/\; \bm{w}_n^\mathsf{H} \bm{u}$ \\
\quad 4. \; $\bm{x}_{n+1} = \bm{x}_n + \alpha_n\,\bm{p}_n$ \\
\quad 5. \; $\bm{r}_{n+1} = \bm{r}_n - \alpha_n\,\bm{u}$, \quad $\bm{q}_{n+1} = \bm{q}_n - \overline{\alpha}_n\,\bm{v}$ \\
\quad 6. \; $s_{n+1} = \bm{q}_{n+1}^\mathsf{H} \bm{r}_{n+1}$ \\
\quad 7. \; If $\mathrm{Re}(\bm{r}_{n+1}^\top \bm{r}_{n+1}) / \|\bm{b}\|^2 < \varepsilon$: \textbf{stop}, return $\bm{x}_{n+1}$. \\
\quad 8. \; $\beta_n = s_{n+1} / s_n$ \\
\quad 9. \; $\bm{p}_{n+1} = \bm{r}_{n+1} + \beta_n\,\bm{p}_n$, \quad $\bm{w}_{n+1} = \bm{q}_{n+1} + \overline{\beta}_n\,\bm{w}_n$ \\
\bottomrule
\end{tabular}
\end{table}

Note: the convergence criterion (step~7) uses the \emph{bilinear} (not sesquilinear) inner product $\bm{r}^\top \bm{r} = \sum_i r_i^2$, following the Fortran MSTM implementation.  The shadow residual $\bm{q} = \overline{\bm{r}}$ ensures that $\bm{q}^\mathsf{H}\bm{r} = \sum_i r_i^2$.

\subsection{GMRES algorithm (alternative solver)}

As an alternative to CBICG, a GMRES (Generalized Minimal Residual) solver (Saad \& Schultz 1986) is available via Krylov.jl (\texttt{solver=:gmres}).  GMRES requires only the forward operator $\mathbf{L}\,\bm{v}$ (no adjoint) and minimises the 2-norm of the residual over a Krylov subspace
\begin{equation}
  \mathcal{K}_k(\mathbf{L}, \bm{r}_0) = \mathrm{span}\{\bm{r}_0,\, \mathbf{L}\,\bm{r}_0,\, \mathbf{L}^2\,\bm{r}_0,\, \ldots,\, \mathbf{L}^{k-1}\,\bm{r}_0\}.
  \label{eq:krylov-subspace}
\end{equation}
The algorithm is summarised in Table~\ref{tab:gmres}.

\begin{table}[htbp]
\caption{GMRES algorithm (Saad \& Schultz 1986).}
\label{tab:gmres}
\centering
\renewcommand{\arraystretch}{1.3}
\begin{tabular}{@{}p{0.95\textwidth}@{}}
\toprule
\textbf{Input:} $\mathbf{L}$, $\bm{b}$, tolerance $\varepsilon$, restart parameter $m$, max iterations $n_\mathrm{max}$. \\
\textbf{Initialise:} $\bm{x}_0 = \bm{b}$, \; $\bm{r}_0 = \bm{b} - \mathbf{L}\,\bm{x}_0$, \; $\beta = \|\bm{r}_0\|_2$, \; $\bm{v}_1 = \bm{r}_0 / \beta$. \\[4pt]
\textbf{For} $j = 1, 2, \ldots, \min(m, n_\mathrm{max})$: \\[2pt]
\quad 1. \; $\bm{w} = \mathbf{L}\,\bm{v}_j$ \hfill (1 forward operator application) \\
\quad 2. \; \textbf{For} $i = 1, \ldots, j$: \; $h_{ij} = \bm{v}_i^\mathsf{H}\,\bm{w}$, \; $\bm{w} \leftarrow \bm{w} - h_{ij}\,\bm{v}_i$ \hfill (modified Gram--Schmidt) \\
\quad 3. \; $h_{j+1,j} = \|\bm{w}\|_2$, \; $\bm{v}_{j+1} = \bm{w} / h_{j+1,j}$ \\
\quad 4. \; Solve $\bm{y}_j = \arg\min_{\bm{y}} \|\beta\,\bm{e}_1 - \bar{\mathbf{H}}_j\,\bm{y}\|_2$ via Givens rotations \\
\quad 5. \; If $\|\beta\,\bm{e}_1 - \bar{\mathbf{H}}_j\,\bm{y}_j\|_2 / \|\bm{b}\|_2 < \varepsilon$: \textbf{stop}. \\[4pt]
\textbf{Solution update:} $\bm{x} = \bm{x}_0 + \mathbf{V}_j\,\bm{y}_j$, where $\mathbf{V}_j = [\bm{v}_1, \ldots, \bm{v}_j]$. \\
If $j = m$ and not converged: restart with $\bm{x}_0 \leftarrow \bm{x}$. \\
\bottomrule
\end{tabular}
\end{table}

Key properties of GMRES compared to CBICG:
\begin{itemize}
  \item 1 operator application per iteration (no adjoint), versus 2 for CBICG.
  \item Monotonically decreasing residual: $\|\bm{r}_k\|_2 \leq \|\bm{r}_{k-1}\|_2$ by construction.
  \item Memory: $O(kn)$ for $k$ Krylov basis vectors of length $n$ (${\sim}\SI{14}{\mega\byte}$ for $k = 30$, $n = 48000$).
  \item The implementation uses $m = 30$ (Krylov.jl \texttt{memory} parameter).
\end{itemize}

\textbf{Performance comparison.}
For the MSTM problem class, GMRES requires approximately twice as many iterations as CBICG to reach the same residual tolerance, because the CBICG adjoint operator (via S-parity, \eqref{eq:adjoint-symmetry}) provides additional directional information.  The total number of $\mathbf{A}$-operator applications is therefore similar (CBICG: $k$ iterations $\times$ 2; GMRES: $\sim 2k$ iterations $\times$ 1), and the wall-clock performance is comparable.  The default solver remains CBICG, consistent with the Fortran MSTM reference implementation.

\subsection{Continuation method for refractive index sweeps}
\label{sec:continuation}

When sweeping the complex refractive index $m$ over a grid while keeping the aggregate geometry fixed, the solution $\bm{x}(m)$ of \eqref{eq:global-system} varies continuously with $m$.  A natural acceleration strategy is to use the converged solution $\bm{x}(m_j)$ at grid point $m_j$ as the initial guess $\bm{x}_0$ for the iterative solver at the neighbouring grid point $m_{j+1}$, instead of the default initial guess $\bm{x}_0 = \bm{b} = \mathbf{T}\,\bm{p}_\mathrm{inc}$.  This is known as a \emph{continuation method} (or warm-start).

When the refractive index changes from $m_j$ to $m_{j+1}$, only the Mie coefficients $a_n, b_n$ (and hence the T-matrix $\mathbf{T}$) change; the translation operator $\mathbf{A}$ depends only on geometry and is unchanged.  For small $|m_{j+1} - m_j|$, the new solution is close to the previous one, so the initial residual $\|\bm{b}_{j+1} - \mathbf{L}_{j+1}\,\bm{x}(m_j)\|$ is much smaller than $\|\bm{b}_{j+1} - \mathbf{L}_{j+1}\,\bm{b}_{j+1}\|$, leading to fewer iterations.

The continuation is applied automatically in the parameter sweep (\texttt{run\_parameter\_sweep}): within each (aggregate, medium condition) group, the solver coefficients $\bm{x}$ from the previous refractive index grid point are passed as the initial guess to the next.  If the equation system size changes (e.g., due to a change in the converged Mie truncation order), the continuation is skipped and the default initial guess is used.

Benchmarks (1000 spheres, FFT mode, $\Delta m \approx 0.05$) show a reduction from 15 to 9--11 iterations per grid point, corresponding to a wall-time speedup of approximately $1.2\times$.  For smaller aggregates ($N_p \sim 100$) where the cold-start solver already converges in 3--4 iterations, the improvement is marginal.

\paragraph{Snake-path grid traversal.}
The effectiveness of the continuation method depends critically on the step size $|m_{j+1} - m_j|$ between consecutive grid points.
For a 2-D grid of $N_r \times N_i$ points spanning the real and imaginary parts of $m$, a naive row-major (raster) traversal sweeps $m_\mathrm{imag}$ from min to max for each value of $m_\mathrm{real}$, then resets to the minimum when advancing to the next $m_\mathrm{real}$ row.
At each row boundary this produces a jump $\Delta m_\mathrm{imag} = m_\mathrm{imag,max} - m_\mathrm{imag,min}$, which can be much larger than the intra-row step.

The \emph{snake-path} (boustrophedon) traversal eliminates this problem by reversing the sweep direction of $m_\mathrm{imag}$ on alternate rows:
\begin{equation*}
  \text{row}~k : \quad m_\mathrm{imag} = \begin{cases}
    m_\mathrm{imag,min} \to m_\mathrm{imag,max} & k~\text{odd}, \\
    m_\mathrm{imag,max} \to m_\mathrm{imag,min} & k~\text{even}.
  \end{cases}
\end{equation*}
This ensures that consecutive grid points always differ by either $\Delta m_\mathrm{real}$ or $\Delta m_\mathrm{imag}$ (one intra-row step), never by the full range.
For the typical CAS sweep grid ($18 \times 26$, $m_\mathrm{real} \in [1.55, 2.4]$, $m_\mathrm{imag} \in [0.15, 1.4]$), the maximum step size drops from $|{\Delta m}| \approx 1.25$ (raster) to $|{\Delta m}| = 0.05$ (snake), a factor of 25 improvement at each of the 17 row boundaries.

\paragraph{GPU-accelerated batched solver with hybrid CPU parallelism.}
An optional CUDA.jl extension provides GPU-accelerated parameter sweeps with hybrid CPU parallelism.
For a fixed (aggregate, medium condition) group, the translation operator $\mathbf{A}$ depends only on geometry and is shared across all refractive index values.
The GPU solver exploits this by processing $B$ RI values simultaneously: it uploads the shared FFT translation data to GPU once, then solves $B$ independent CBICG systems $(\mathbf{I} - \mathbf{T}_b \mathbf{A})\bm{x}_b = \mathbf{T}_b\bm{p}_\mathrm{inc}$ in parallel, where only $\mathbf{T}_b$ (the Mie T-matrix) varies per RI.
$\mathbf{T}_b$ holds the Mie coefficients of every sphere at refractive index $b$, so spheres of unequal radii are treated exactly.
Versions before v1.1.0 used the coefficients of the first sphere for all spheres in the iteration operator, which was correct for equal spheres only.

For large aggregates ($N_p \gtrsim 500$), the CPU phases of each batch (preparation of RHS vectors and post-processing of solutions) dominate wall time and leave the GPU idle.
This bottleneck is addressed by three complementary measures:
\begin{enumerate}
  \item \textbf{Intra-group CPU parallelism}: batch preparation and post-processing are parallelised over the $B$ batch items using \texttt{Threads.@threads}, reducing the CPU-phase time by a factor of $C$ (number of Julia threads).
  \item \textbf{Hybrid GPU+CPU dispatch}: groups with $N_p \leq$ \texttt{gpu\_np\_threshold} (default 500) are processed by CPU threads running \emph{concurrently} with the GPU task via \texttt{Threads.@spawn}, so no CPU core sits idle while the GPU handles a large-aggregate batch.
  \item \textbf{Batched common-origin translation} (v0.4.1): the common-origin translation matrices $\mathbf{J}_{0i}$ (Eq.~\eqref{eq:common-origin}) depend only on geometry, not on the refractive index.
  They are computed once per group (not once per RI value) and applied to all $B$ solutions simultaneously via a BLAS GEMM,
  \[
    (\mathbf{J}_{0i}\,\bm{a}^{(i)})_\text{batch} = \mathbf{J}_{0i} \times \bm{A}_\text{batch},
  \]
  where $\bm{A}_\text{batch}$ has shape $(N_\text{hnb}, 2B)$.
  This reduces the calls of \texttt{compute\_translation\_matrix} from $N_S \times n_\text{RI}$ to $N_S$ per group.
\end{enumerate}
Running with \texttt{julia -t auto} activates all mechanisms.
Measured with v0.4.1 on an RTX~A6000 with a 24-core CPU (CAS sweep, $N_p = 350$--$2800$, FP32 mode), the hybrid implementation achieved ${\sim}6$ jobs/s versus ${\sim}0.8$ jobs/s without threading, a factor of ${\sim}7$.
With v1.1.0 a sweep of \num{2.02e6} jobs took \SI{30}{\hour} on the same machine ($N_p = 350$--$2800$ in steps of 50, 13 fractal dimensions with 10 aggregates each, three wavelengths with 90 to 121 refractive indices per aggregate, uniform order 3, FP32 mode).
Three GPU processes with 8 threads each computed the groups with $N_p \geq 550$ at 17 jobs/s in total.
One CPU process with 6 threads computed the groups with $N_p \leq 500$ at 5.4 jobs/s.
The mean number of iterations per job was 7 to 10 on the GPU.

A Float32 precision mode is available (\texttt{gpu\_float32=true}): Mie coefficients and the RHS are computed in Float64 on CPU, while the iterative solver runs in ComplexF32 on GPU.
The convergence tolerance is automatically relaxed to $10^{-5}$.
Validation against the Fortran reference for $N_S = 1000$ spheres across the full CAS RI range shows that FP32 errors are dominated by the FFT translation approximation, not by the reduced floating-point precision.

\subsection{Incident plane wave coefficients}
\label{sec:incident}

For $z$-axis incidence ($\beta = 0$, $\alpha = 0$), only the $m = \pm 1$ multipole components are non-zero due to the azimuthal symmetry of a circularly decomposed plane wave.
In the lr\_tran basis, the incident field coefficients for sphere $i$ are:
\begin{equation}
  p_{mn,p=1,q=1}^{(i)} = i^{\,n+1}\sqrt{\frac{2n+1}{2}}\;\delta_{m,+1}\;e^{iz_i}, \qquad
  p_{mn,p=2,q=1}^{(i)} = -i^{\,n+1}\sqrt{\frac{2n+1}{2}}\;\delta_{m,-1}\;e^{iz_i}
  \label{eq:pinc-q1}
\end{equation}
for the first incident polarisation ($q = 1$), and
\begin{equation}
  p_{mn,p=1,q=2}^{(i)} = -i^{\,n+2}\sqrt{\frac{2n+1}{2}}\;\delta_{m,+1}\;e^{iz_i}, \qquad
  p_{mn,p=2,q=2}^{(i)} = -i^{\,n+2}\sqrt{\frac{2n+1}{2}}\;\delta_{m,-1}\;e^{iz_i}
  \label{eq:pinc-q2}
\end{equation}
for $q = 2$.
The phase factor $e^{iz_i}$ accounts for the spatial position of sphere $i$ along the propagation direction.

%% ============================================================
\section{Output Quantities}
\label{sec:output}

\subsection{Translation to a common origin}

After solving for the per-sphere scattered-field coefficients $\bm{a}^{(i)}$, all coefficients are re-expanded about a common origin $\bm{r}_0$ (typically the aggregate centroid) using regular (Bessel $j_n$) translation:
\begin{equation}
  \bm{a}_0 = \sum_{i=1}^{N_S} \mathbf{J}_{0i}\,\bm{a}^{(i)}
  \label{eq:common-origin}
\end{equation}
where $\mathbf{J}_{0i}$ is the regular translation matrix from sphere $i$ to the common origin.
The translation order for each sphere is determined by a convergence test on the translation addition theorem (Fortran routine \texttt{tranordertest}), ensuring that the re-expansion captures the scattered field to a specified accuracy.

The common-origin lr\_tran coefficients are then transformed from the $(p=1, p=2)$ basis to mode-indexed form:
\begin{equation}
  a_{mn}^{(\text{mode}\,1)} = a_{mn}^{(p=1)} + a_{mn}^{(p=2)}, \qquad
  a_{mn}^{(\text{mode}\,2)} = a_{mn}^{(p=1)} - a_{mn}^{(p=2)}.
  \label{eq:mode-transform}
\end{equation}
These mode-1 and mode-2 coefficients correspond to the standard TE and TM VSWF expansion of the total scattered field about the common origin (Mackowski 1994, \S 3).

\subsection{Scattering amplitude matrix (BH83 convention)}

The far-field scattered wave is described by the $2 \times 2$ amplitude matrix (BH83, Eq.~3.12):
\begin{equation}
  \begin{pmatrix} E_\parallel^\mathrm{sca} \\ E_\perp^\mathrm{sca} \end{pmatrix}
  = \frac{e^{ikr}}{-ikr}
  \begin{pmatrix} S_2 & S_3 \\ S_4 & S_1 \end{pmatrix}
  \begin{pmatrix} E_\parallel^\mathrm{inc} \\ E_\perp^\mathrm{inc} \end{pmatrix}
  \label{eq:bh83-amplitude}
\end{equation}
where $S_1, S_2, S_3, S_4$ are dimensionless complex functions of the scattering angles $(\theta, \phi)$.

For the forward ($\theta = 0^\circ$) and backward ($\theta = 180^\circ$) directions only the $m = \pm 1$ terms of the common-origin expansion survive.
The amplitudes are computed from the mode coefficients as
\begin{equation}
  S_j = -2 \sum_{n=1}^{N_\mathrm{max}} \sum_{m = \pm 1} \sum_{p = 1}^{2}
  F_j(n, m, p)\;\tau_{m,n,p}(\theta)
  \label{eq:amplitude-sum}
\end{equation}
where $F_j$ selects the appropriate combination of mode coefficients and incident polarisation (mode-1 from $q=1$, mode-2 from $q=1$ for $S_2$, $S_4$; similarly with sign change from $q=2$ for $S_1$, $S_3$), and $\tau_{m,n,p}(\theta)$ are angular functions derived from the Wigner $d$-matrix at $\cos\theta$.

At $\theta = 0^\circ$ the angular functions simplify to (Mackowski \& Mishchenko 2011, \texttt{taufunc} routine):

\begin{center}
\begin{tabular}{@{} ccc @{}}
  \toprule
  $m$ & $p$ & $\tau_{m,n,p}(0^\circ)$ \\
  \midrule
  $+1$ & 1 & $-f_n$ \\
  $+1$ & 2 & $-f_n$ \\
  $-1$ & 1 & $+f_n$ \\
  $-1$ & 2 & $-f_n$ \\
  \bottomrule
\end{tabular}
\end{center}
where $f_n = \frac{1}{4}\sqrt{(2n+1)/2}$.

At $\theta = 180^\circ$ an additional factor of $(-1)^n$ appears, and the signs of the $p = 1$ vs $p = 2$ contributions differ:

\begin{center}
\begin{tabular}{@{} ccc @{}}
  \toprule
  $m$ & $p$ & $\tau_{m,n,p}(180^\circ)$ \\
  \midrule
  $+1$ & 1 & $-(-1)^n f_n$ \\
  $+1$ & 2 & $+(-1)^n f_n$ \\
  $-1$ & 1 & $+(-1)^n f_n$ \\
  $-1$ & 2 & $+(-1)^n f_n$ \\
  \bottomrule
\end{tabular}
\end{center}

\textbf{Note on $S_3$ and $S_4$:}
For a single sphere or a symmetric dimer with the symmetry axis along $z$, $S_3 = S_4 = 0$ by symmetry.
For general aggregates lacking such symmetry, all four elements are non-zero.

\subsection{Cross-section efficiencies}

The dimensionless efficiency factors are defined as $Q = C / (\pi\,a_\mathrm{eff}^2)$ where $C$ is the corresponding cross section, and computed for unpolarised incidence by averaging over the two incident polarisations $q = 1, 2$.

\textbf{Extinction efficiency} (optical theorem, Mackowski 1994, Eq.~12):
\begin{equation}
  Q_\mathrm{ext} = -\frac{2}{x_\mathrm{eff}^2}
  \sum_{q=1}^{2} \mathrm{Re}\!\left(\bm{a}_q^\dagger \,\bm{p}_{\mathrm{inc},q}\right)
  \label{eq:qext}
\end{equation}
where the inner product $\bm{a}_q^\dagger \bm{p}_{\mathrm{inc},q}$ is evaluated in the per-sphere basis (before common-origin translation).
This is the multi-sphere generalisation of the forward-scattering (optical) theorem: the extinction is determined entirely by the projection of the scattered-field coefficients onto the incident-field pattern.
The factor of 2 (instead of 4) in the prefactor arises from averaging the cross section over the two incident orthogonal polarisation states ($q=1, 2$).

\textbf{Scattering efficiency} (from the coherent common-origin expansion, Mackowski 1994, Eq.~13):
\begin{equation}
  Q_\mathrm{sca} = \frac{1}{x_\mathrm{eff}^2}
  \sum_{q=1}^{2} \left(\|\bm{a}_{0,q}^{(\text{mode}\,1)}\|^2 + \|\bm{a}_{0,q}^{(\text{mode}\,2)}\|^2 \right)
  \label{eq:qsca}
\end{equation}
where the norm is the sum of squared magnitudes of all common-origin expansion coefficients for incident polarisation $q$.
This requires the common-origin re-expansion (Sec.~\ref{sec:output}, first subsection), which coherently accounts for the relative phases of all spheres.

\textbf{Absorption efficiency:}
\begin{equation}
  Q_\mathrm{abs} = Q_\mathrm{ext} - Q_\mathrm{sca}
  \label{eq:qabs}
\end{equation}
following from energy conservation.

\subsection{MI02 scattering amplitudes}

The MI02 (Mishchenko, Travis \& Lacis 2002) amplitude matrix $\mathbf{S}^\mathrm{MI02}$ has dimension of length and is related to the dimensionless BH83 amplitudes by:
\begin{equation}
  S_{11}^\mathrm{MI02} = \frac{S_2^\mathrm{BH83}}{-ik}, \quad
  S_{22}^\mathrm{MI02} = \frac{S_1^\mathrm{BH83}}{-ik}, \quad
  S_{12}^\mathrm{MI02} = \frac{S_3^\mathrm{BH83}}{ik}, \quad
  S_{21}^\mathrm{MI02} = \frac{S_4^\mathrm{BH83}}{ik}
  \label{eq:mi02-conversion}
\end{equation}
where $k = 2\pi\,n_\mathrm{med}/\lambda_0$ is the (dimensional) wavenumber in the medium.
Note the index mapping: the MI02 $S_{11}$ element corresponds to BH83 $S_2$ (not $S_1$), and vice versa.

\subsection{CAS-v2 observable amplitudes}

The Complex Amplitude Sensing (CAS-v2) protocol (Moteki \& Adachi 2024) directly measures three complex amplitudes derived from the MI02 matrix elements:
\begin{align}
  S_s^\mathrm{fwd} &= S_{11}^\mathrm{MI02} + i\,S_{12}^\mathrm{MI02} \label{eq:cas-ss} \\
  S_p^\mathrm{fwd} &= S_{22}^\mathrm{MI02} - i\,S_{21}^\mathrm{MI02} \label{eq:cas-sp} \\
  S^\mathrm{bak}   &= \frac{-S_{11}^\mathrm{MI02,bwd} + S_{22}^\mathrm{MI02,bwd}}{2} \label{eq:cas-sbak}
\end{align}
where $S_s^\mathrm{fwd}$ and $S_p^\mathrm{fwd}$ are the forward scattering amplitudes for $s$- and $p$-polarised incident light, respectively, and $S^\mathrm{bak}$ is the depolarisation-sensitive backward scattering amplitude.
Equation~\eqref{eq:cas-sbak} is the corrected Eq.~(28) of Moteki \& Adachi (2024), given in the erratum to that paper.
The printed equation has $1/\sqrt{2}$ in place of $1/2$.

%% ============================================================
\section{Cluster T-Matrix and Particle Orientation}
\label{sec:cluster-tmatrix}

The solution of Eq.~\eqref{eq:global-system} holds for one orientation of the aggregate relative to the incident wave.
Another orientation can be obtained by rotating the sphere positions and solving again.
For many orientations of one aggregate it is cheaper to build the T-matrix of the whole aggregate once.
The functions \texttt{build\_cluster\_tmatrix} and \texttt{cluster\_forward\_amplitudes} implement this path.

\subsection{Cluster T-matrix}

The cluster T-matrix $\mathbf{T}^\mathrm{cl}$ maps the coefficients $\bm{p}_0$ of a field that is regular about the common origin $\bm{r}_0$ to the coefficients $\bm{a}_0$ of the scattered field about the same origin:
\begin{equation}
  \bm{a}_0 = \mathbf{T}^\mathrm{cl}\,\bm{p}_0.
  \label{eq:cluster-tmatrix}
\end{equation}
Both vectors are in the lr\_tran basis with the combined index $(mn, p)$.
The origin $\bm{r}_0$ is the centroid of the sphere centres.
The order $N_0$ of the expansion about the origin is the largest translation order of Sec.~\ref{sec:output} over all spheres, so $\mathbf{T}^\mathrm{cl}$ is a square matrix of size $K = 2N_0(N_0+2)$.
The matrix depends on the geometry, the size and $m_\mathrm{rel}$.
It does not depend on the orientation.

Column $\nu$ of $\mathbf{T}^\mathrm{cl}$ is the response to a unit regular VSWF of index $\nu$ at the origin.
The mode is translated to each sphere by the regular translation matrix $\mathbf{J}_{i0}$,
\begin{equation}
  \bm{p}_\nu^{(i)} = \mathbf{J}_{i0}\,\bm{e}_\nu,
  \label{eq:cluster-rhs}
\end{equation}
where $\bm{e}_\nu$ is the unit vector of mode $\nu$.
The interaction equation~\eqref{eq:global-system} is solved with this right-hand side, and the solution is merged to the origin by Eq.~\eqref{eq:common-origin}:
\begin{equation}
  \mathbf{T}^\mathrm{cl}\,\bm{e}_\nu = \sum_{i=1}^{N_S} \mathbf{J}_{0i}\,\bm{a}_\nu^{(i)}.
  \label{eq:cluster-column}
\end{equation}
All $K$ right-hand sides share the operator $\mathbf{I} - \mathbf{T}\mathbf{A}$.
The present implementation uses the direct translation operator for these solves, not the FFT operator.

\subsection{Rotation of the expansion coefficients}

The orientation of the particle is given by the intrinsic ZYZ Euler angles $\Omega = (\alpha, \beta, \gamma)$.
A rotation mixes VSWFs of the same order $n$ only, and it acts in the same way on the two polarisation blocks.
Within order $n$ the coefficients transform with the Wigner D-matrix (see, e.g., MI02, Appendix~B)
\begin{equation}
  D^n_{m'm}(\alpha, \beta, \gamma) = e^{-i m' \alpha}\, d^n_{m'm}(\beta)\, e^{-i m \gamma},
  \label{eq:wigner-D}
\end{equation}
where $d^n_{m'm}(\beta)$ is the real Wigner d-matrix.
Let $\mathbf{D}(\Omega)$ denote the block-diagonal matrix of these blocks for $n = 1, \dots, N_0$ and both values of $p$.
It is unitary.

Each element $d^n_{m'm}(\beta)$ is computed by an upward recurrence in $n$.
The recurrence starts at the lowest order $n_\mathrm{min} = \max(|m'|, |m|)$.
At this order Wigner's formula
\begin{multline}
  d^n_{m'm}(\beta) = \sum_k (-1)^{m'-m+k}\,
  \frac{\sqrt{(n+m)!\,(n-m)!\,(n+m')!\,(n-m')!}}{(n+m-k)!\,k!\,(m'-m+k)!\,(n-m'-k)!} \\
  \times \left(\cos\frac{\beta}{2}\right)^{2n+m-m'-2k} \left(\sin\frac{\beta}{2}\right)^{m'-m+2k}
  \label{eq:wigner-formula}
\end{multline}
has a single term, so the seed is free of cancellation.
With $x = \cos\beta$ and $d^{\,n_\mathrm{min}-1}_{m'm} = 0$ the higher orders follow from
\begin{multline}
  l\,\sqrt{\left[(l+1)^2 - m'^2\right]\left[(l+1)^2 - m^2\right]}\; d^{\,l+1}_{m'm}
  = (2l+1)\left[l(l+1)\,x - m' m\right] d^{\,l}_{m'm} \\
  - (l+1)\,\sqrt{\left(l^2 - m'^2\right)\left(l^2 - m^2\right)}\; d^{\,l-1}_{m'm}.
  \label{eq:wigner-recurrence}
\end{multline}
For $m' = m = 0$ the element is the Legendre polynomial $P_n(x)$ and its own recurrence is used.

\subsection{Forward amplitudes at a given orientation}

The incident wave stays along $+z$ in the laboratory frame and the particle is rotated by $\Omega$.
The scattered coefficients about the origin are
\begin{equation}
  \bm{a}_0(\Omega) = \mathbf{D}(\Omega)\,\mathbf{T}^\mathrm{cl}\,\mathbf{D}(\Omega)^\dagger\,\bm{p}_{0,z},
  \label{eq:cluster-orientation}
\end{equation}
where $\bm{p}_{0,z}$ holds the plane-wave coefficients of Sec.~\ref{sec:incident} at the origin.
They carry the phase factor $e^{i z_0}$, where $z_0$ is the dimensionless $z$ coordinate of $\bm{r}_0$.
The phase reference is then the laboratory plane $z = 0$, as in the fixed-orientation solution.
Equation~\eqref{eq:cluster-orientation} is evaluated from right to left as two rotations of a vector and one product of a matrix and a vector.
The cost per orientation is $O(K^2)$.
The forward amplitudes $S_1$ to $S_4$ follow from $\bm{a}_0(\Omega)$ by Eq.~\eqref{eq:mode-transform} and the amplitude sums of Sec.~\ref{sec:output}.
$\Omega = (0, 0, 0)$ reproduces the fixed-orientation result.

The implementation was checked in two ways.
At $\Omega = (0,0,0)$ the amplitudes agree with the direct solution of Eq.~\eqref{eq:global-system}.
At other orientations they agree with the direct solution for the geometry rotated about its centroid by the same angles.
The present functions return the forward amplitudes only.

%% ============================================================
\section{Computational Flow Summary}
\label{sec:flow}

\begin{center}
\footnotesize
\begin{tabular}{@{} rlll @{}}
  \toprule
  Step & Operation & Description & Source file \\
  \midrule
  1 & Non-dimensionalise & $x_i = k \cdot a_i$, $\bm{r}_i \to k \cdot \bm{r}_i$, $m_\mathrm{rel} = m_\mathrm{sphere}/n_\mathrm{med}$ & \texttt{ParameterSweep.jl} \\
  2 & Mie coefficients & $a_n, b_n$ for each sphere & \texttt{MieCoefficients.jl} \\
  3 & Build T-matrix & lr\_tran basis, per sphere & \texttt{TMatrixSolver.jl} \\
  4 & Incident field & $\bm{p}_\mathrm{inc}$ ($m = \pm 1$ terms) & \texttt{TMatrixSolver.jl} \\
  5 & CBICG solve & $(\mathbf{I} - \mathbf{T}\mathbf{A})\bm{x} = \mathbf{T}\bm{p}_\mathrm{inc}$; warm-start from previous $m$ (\S\ref{sec:continuation}) & \texttt{TMatrixSolver.jl} \\
  6 & Common origin & Translate coefficients & \texttt{ScatteringAmplitude.jl} \\
  7 & Amplitudes & $S_1$--$S_4$ at $\theta = 0^\circ, 180^\circ$ & \texttt{ScatteringAmplitude.jl} \\
  8 & Efficiencies & $Q_\mathrm{ext}$, $Q_\mathrm{sca}$, $Q_\mathrm{abs}$ & \texttt{ScatteringAmplitude.jl} \\
  9 & MI02 \& CAS & Convert amplitudes & \texttt{ParameterSweep.jl} \\
  \bottomrule
\end{tabular}
\end{center}

%% ============================================================
\section{Implementation Mapping}
\label{sec:mapping}

\begin{center}
\footnotesize
\let\origunderscore\_
\renewcommand{\_}{\origunderscore\allowbreak}
\begin{tabular}{@{} >{\raggedright\arraybackslash}p{3.5cm} >{\raggedright\arraybackslash}p{2.9cm} >{\raggedright\arraybackslash}p{4.7cm} >{\raggedright\arraybackslash}p{4.4cm} @{}}
  \toprule
  Theory & Equation(s) & Implementation & File \\
  \midrule
  Wavenumber, non-dim. & \eqref{eq:wavenumber}--\eqref{eq:a-eff} & \texttt{compute\_scattering}, \texttt{run\_parameter\_sweep} & \texttt{ScatteringAmplitude.jl}, \texttt{ParameterSweep.jl} \\
  Mie coefficients $a_n, b_n$ & \eqref{eq:mie-an}--\eqref{eq:nstop} & \texttt{compute\_mie\_coefficients} & \texttt{MieCoefficients.jl} \\
  lr\_tran T-matrix & \eqref{eq:tmatrix-lr} & \texttt{\_precompute\_T\_values}, \texttt{\_apply\_T!} & \texttt{TMatrixSolver.jl} \\
  Translation $\mathbf{H}_{ij}$ & \eqref{eq:axial-translation}--\eqref{eq:lr-translation} & \texttt{compute\_translation\_matrix} & \texttt{TranslationCoefs.jl} \\
  FFT translation & \eqref{eq:fft-convolution} & \texttt{apply\_A\_fft!}, \texttt{init\_fft\_grid} & \texttt{FFTTranslation.jl} \\
  CBICG solver & \eqref{eq:global-system}, \eqref{eq:adjoint-symmetry}, Table~\ref{tab:cbicg} & \texttt{\_solve\_bicg} & \texttt{TMatrixSolver.jl} \\
  GMRES solver (optional) & \eqref{eq:global-system}, \eqref{eq:krylov-subspace}, Table~\ref{tab:gmres} & \texttt{gmres} (Krylov.jl) & \texttt{TMatrixSolver.jl} \\
  RI continuation (\S\ref{sec:continuation}) & --- & \texttt{initial\_amn} kwarg & \texttt{TMatrixSolver.jl}, \texttt{ParameterSweep.jl} \\
  Incident coefficients & \eqref{eq:pinc-q1}--\eqref{eq:pinc-q2} & \texttt{\_genplanewavecoef\_z0} & \texttt{TMatrixSolver.jl} \\
  Common-origin expansion & \eqref{eq:common-origin}--\eqref{eq:mode-transform} & \texttt{\_merge\_to\_origin!} & \texttt{ScatteringAmplitude.jl} \\
  BH83 amplitudes $S_1$--$S_4$ & \eqref{eq:bh83-amplitude}--\eqref{eq:amplitude-sum} & \texttt{\_amplitude\_from\_mode\_coefs} & \texttt{ScatteringAmplitude.jl} \\
  Cross-section efficiencies & \eqref{eq:qext}--\eqref{eq:qabs} & \texttt{compute\_scattering} & \texttt{ScatteringAmplitude.jl} \\
  MI02 conversion & \eqref{eq:mi02-conversion} & \texttt{run\_parameter\_sweep} (inline) & \texttt{ParameterSweep.jl} \\
  CAS-v2 observables & \eqref{eq:cas-ss}--\eqref{eq:cas-sbak} & \texttt{read\_results\_h5\_for\_cas.jl} & \texttt{scripts/} \\
  Aggregate I/O & --- & \texttt{read\_aggregate\_file}, \texttt{read\_aggregate\_catalog} & \texttt{AggregateIO.jl} \\
  Contact order (\S\ref{sec:contact-order}) & \eqref{eq:contact-distance}--\eqref{eq:contact-error} & \texttt{contact\_order\_scan}, \texttt{contact\_truncation\_order} & \texttt{TruncationOrder.jl} \\
  Far order of the grid & \eqref{eq:near-pair} & \texttt{init\_fft\_grid} (\texttt{far\_nois}, \texttt{near\_factor}) & \texttt{FFTTranslation.jl} \\
  Cluster T-matrix & \eqref{eq:cluster-tmatrix}--\eqref{eq:cluster-column} & \texttt{build\_cluster\_tmatrix} & \texttt{ClusterTMatrix.jl} \\
  Wigner rotation & \eqref{eq:wigner-D}--\eqref{eq:wigner-recurrence} & \texttt{wigner\_d}, \texttt{wigner\_D} & \texttt{WignerRotation.jl} \\
  Amplitudes at an orientation & \eqref{eq:cluster-orientation} & \texttt{cluster\_forward\_amplitudes} & \texttt{ClusterTMatrix.jl} \\
  \midrule
  \multicolumn{4}{l}{\textit{GPU extension (optional, via \texttt{using CUDA})}} \\
  GPU batched BiCG & \eqref{eq:global-system}, Table~\ref{tab:cbicg} & \texttt{gpu\_solve\_bicg\_batch!} & \texttt{ext/.../gpu\_batch\_solver.jl} \\
  GPU FFT translation & \eqref{eq:fft-convolution} & \texttt{gpu\_apply\_A\_batch!} & \texttt{ext/.../gpu\_fft\_translation.jl} \\
  GPU T-matrix / S-parity & \eqref{eq:tmatrix-lr}, \eqref{eq:adjoint-symmetry} & \texttt{gpu\_apply\_T\_batch!} & \texttt{ext/.../gpu\_kernels.jl} \\
  \bottomrule
\end{tabular}
\end{center}

%% ============================================================
\section{References}
\label{sec:references}

\begin{enumerate}[nosep]
  \item Bohren, C.F.\ \& Huffman, D.R.\ (1983). \textit{Absorption and Scattering of Light by Small Particles}. Wiley. [BH83]
  \item Mackowski, D.W.\ (1991). Analysis of radiative scattering for multiple sphere configurations. \textit{Proc.\ R.\ Soc.\ Lond.\ A}, 433, 599--614.
  \item Mackowski, D.W.\ (1994). Calculation of total cross sections of multiple-sphere clusters. \textit{JOSA A}, 11(11), 2851--2861.
  \item Mackowski, D.W.\ (1996). Calculation of the total scattering cross section and the total absorbed energy for the multiple sphere clusters. \textit{JOSA A}, 13(11), 2266--2278.
  \item Mackowski, D.W.\ (2014). A general superposition solution for electromagnetic scattering by multiple spherical domains of optically active media. \textit{JQSRT}, 133, 264--270.
  \item Mackowski, D.W.\ \& Kolokolova, L.\ (2022). Application of the multiple sphere superposition solution to large-scale systems of spheres via an accelerated algorithm. \textit{JQSRT}, 280, 108075.
  \item Mackowski, D.W.\ \& Mishchenko, M.I.\ (2011). A multiple sphere T-matrix Fortran code for use on parallel computer clusters. \textit{JQSRT}, 112(13), 2182--2192.
  \item Mishchenko, M.I., Travis, L.D.\ \& Lacis, A.A.\ (2002). \textit{Scattering, Absorption, and Emission of Light by Small Particles}. Cambridge University Press. [MI02]
  \item Moteki, N.\ \& Adachi, K.\ (2024). CAS-v2 protocol. \textit{Optics Express}, 32(21), 36500--36522.
  \item Saad, Y.\ \& Schultz, M.H.\ (1986). GMRES: A generalized minimal residual algorithm for solving nonsymmetric linear systems. \textit{SIAM J.\ Sci.\ Stat.\ Comput.}, 7(3), 856--869.
  \item Wiscombe, W.J.\ (1980). Improved Mie scattering algorithms. \textit{Applied Optics}, 19(9), 1505--1509.
  \item Xu, Y.-l.\ (1996). Calculation of the addition coefficients in electromagnetic multisphere-scattering theory. \textit{J.\ Comput.\ Phys.}, 127(2), 285--298.
\end{enumerate}

\bigskip
\noindent\textbf{Acknowledgment.}\quad
This document was prepared with the assistance of Claude (Anthropic). The author assumes full responsibility for the content.

\end{document}
